\chapter{The Lean 4 companion: verified statements and explicit limits}

\section{Current six-module verification}

The current companion contains 115 selected theorem declarations in six
modules.
\begin{center}
\begin{tabular}{lr}
\toprule
Module & Declarations\\\midrule
\texttt{BridgeTheorems} & 25\\
\texttt{QuantumInterfaces} & 7\\
\texttt{DiagramGrammar} & 16\\
\texttt{FluidInterfaces} & 34\\
\texttt{ReductionContracts} & 16\\
\texttt{TransportContracts} & 17\\\bottomrule
\end{tabular}
\end{center}
The final verification records a fresh compiler check, exact printed types,
source hashes and axiom dependencies for every declaration, with no warnings
or unexpected axioms. Its authoritative record is
\texttt{lean/final-verification/verification.json} and the six module logs.
The formal type map identifies 68 declarations with precise partial
components of 25 manuscript environments; no whole composite environment is
labelled completely formalized. This verification covers the selected local
interfaces, not the full upstream continuum proof or cryptographic security.

The following section preserves the original 25-declaration module's exact
verification record, which remains a component of the current package.

\section{Original module verification result}

The file \texttt{lean/BridgeTheorems.lean} contains 25 named theorem
declarations. It was compiled successfully with Lean 4.24.0 on Linux, using
mathlib revision
\begin{center}\small\texttt{f897ebcf72cd16f89ab4577d0c826cd14afaafc7}.\end{center}
The compiler exited with status zero. Every declaration has a corresponding
\texttt{\#print axioms} command in the audit block. The resulting dependencies
are subsets of the standard set
\[
 \{\texttt{propext},\ \texttt{Classical.choice},\ \texttt{Quot.sound}\}.
\]
No declaration in this companion uses \texttt{sorry}, \texttt{sorryAx}, or a
new project axiom. The actual compiler output is preserved in
\texttt{lean/verification.log}; the toolchain, source hash, declaration names
and axiom lists are preserved in \texttt{lean/verification.json}. Those files
describe a completed local check, not a proposed future build.

The source hash of the verified Lean file is
\begin{center}\scriptsize
\texttt{a85b897824b333fd8e7962bfb7959875e1c7e5d7870e4694b02e345855b895ba}.
\end{center}
If the file is edited, that hash and the build evidence must be regenerated.
The pinned toolchain and dependency commit make the verification target
unambiguous. They do not by themselves prove that a later modified file has
been checked.

\section{Why this is selected formalization}

The present development formalizes an elementary mathematical spine through
the monograph. It does not formalize every theorem cited or proved in the
paper. In particular it does not contain a formal proof of the full
Buchsbaum--Eisenbud acyclicity criterion, the Atiyah--Singer index theorem,
the correctness and complexity of Shor's physical circuit model, a worst-case
LWE security reduction, the inherited arithmetic determinant conjectures, or
the Yang--Mills construction. It proves neither GRH nor a Yang--Mills mass
gap. Stating this scope is necessary to interpret what the compiler has
actually certified.

There are three different proof layers in the deliverable. First, the paper
gives complete elementary proofs of many finite and algebraic results.
Second, it states deep external theorems with citations and their relevant
hypotheses. Third, the current companion checks 115 selected formal declarations
across the six modules listed above. The original four modules contribute 82;
the reduction and transport modules add 16 and 17, respectively.
Appendix~\ref{ch:fluid-verification} identifies the earlier fluid-era checks
and distinguishes them from the current aggregate verification.
These layers overlap, but they are not identical. A citation to a formal
upstream project is also different from compiling that project during this
work. The Yang--Mills dependency manifest records the frozen external source
and assessment, and explicitly states that the full external development
was not rebuilt here.

\section{Linear complexes and rank}

The first group of declarations works with arbitrary linear maps
$A:V_0\to V_1$ and $B:V_1\to V_2$ over a field. The zero-composition
hypothesis is an equality of linear maps, not a numerical tolerance:
\[
 B\circ A=0.
\]
The declaration \path{complex_range_le_kernel} proves
$\im A\subseteq\ker B$. Under finite dimensionality of the middle space,
\path{complex_rank_bound} proves the adjacent-rank inequality, and
\path{middle_exact_iff_rank_sum} proves that equality is equivalent to
middle exactness. These statements are useful formal interfaces because
they distinguish the complex condition from the stronger exactness
condition.

The Euler-rank identity is stated over signed integers. This avoids a common
formalization trap: subtraction on natural numbers is truncated, whereas
the ordinary algebraic cancellation of alternating dimensions assumes
signed arithmetic or hypotheses ensuring every subtraction is valid. The
geometric interpretation of the ranks as dimensions of actual cohomology
groups is supplied by the preceding linear-algebra results and the paper.

\section{The cone and its tangent calculation}

The definition \texttt{cone} is the polynomial
$x_1y_1+x_2y_2$ over a commutative ring. Its perturbation theorem is an exact
polynomial identity with a quadratic remainder. Given a point on the cone
and a direction satisfying the linearized equation, the second theorem
identifies the residual as $t^2(u_1v_1+u_2v_2)$. It does not incorrectly
assert that every tangent vector gives a straight line lying on the
variety.

The sum-and-difference identity supplies the algebra behind the real
torus-link parameterization. It is stated without division, so it is valid
over every commutative ring. The geometric step that divides by $\sqrt2$ and
identifies circles takes place over $\R$ in the paper. The Lean theorem
checks the polynomial part; it does not claim a formal homeomorphism of the
link with a torus. Similarly, the two-generator Koszul declaration proves
zero composition only. Regular-sequence hypotheses are needed for an
exactness theorem and are not silently inferred from commutativity.

\section{Hodge positivity and the characteristic polynomial}

The generic matrix theorem proves that
\[
 AA^*+B^*B
\]
is positive semidefinite for finite real matrices. This statement does not
require $BA=0$. That is mathematically appropriate: positivity of the sum of
two squares is weaker than its interpretation as the Laplacian of a
complex. The paper uses the complex condition when identifying its kernel
with cohomology.

The coordinate calculation uses $A=(x_1,x_2)^T$ and
$B=(y_1,y_2)$ and proves the quadratic-form identity
\[
 z^T\Delta_1z=(x_1z_1+x_2z_2)^2+(y_1z_1+y_2z_2)^2.
\]
The energy is nonnegative, and vanishes exactly when both displayed linear
forms vanish. The characteristic determinant is checked before imposing
the complex equation:
\[
 \det(tI-\Delta_1)
 =(t-\|x\|^2)(t-\|y\|^2)-(x\cdot y)^2.
\]
On the cone $x\cdot y=0$, it factors into the two squared-length factors.
Retaining the defect term in the general identity helps prevent accidental
use of the factorization off the variety.

The final Hodge declarations check the diagonal family and the existence of
arbitrarily small energy along a fixed unit direction. This is a formal
finite-dimensional obstruction to inferring a uniform coercivity constant
from exactness alone. It is not a formal theorem about the spectrum of an
interacting quantum-field Hamiltonian.

\section{Integral shears, cyclic kernels and noisy arithmetic}

The shear declarations compute the determinant, inverse, image of a basis
vector, and its squared norm. Together they make precise the distinction
between an invertible integral map and a map with bounded Euclidean
distortion. They do not formalize a complete lattice-reduction algorithm or
its asymptotic runtime.

The cyclic-kernel declarations work in \texttt{ZMod q}. They prove the
algebraic relation underlying
$aP+b(dP)=(a+bd)P$ and the collision equation. The formal file does not
construct a quantum state, a unitary Fourier transform, a reversible
elliptic-curve circuit, or a success-probability measure. Those require a
substantially larger formal specification. The Fourier proof in the paper
and the finite numerical experiment have their own clearly stated status.

The noise-cancellation declaration is a matrix identity over a commutative
ring, using compatible finite index types. It proves that the public linear
terms cancel and isolates the remaining error. The binary-noise bound is a
real inequality used in an elementary correctness argument. Neither proves
computational indistinguishability, chosen-ciphertext security, or the
security of ML-KEM. Correctness algebra and security reductions answer
different questions.

\section{Reading an axiom audit}

The appearance of \texttt{Classical.choice} in a declaration does not mean
that an arbitrary new cryptographic or analytic hypothesis has been added.
It identifies use of Lean's standard classical foundation. Likewise
\texttt{propext} and \texttt{Quot.sound} are standard logical principles in
the mathlib environment. What would materially change this audit is a new
project axiom asserting a difficult target, or an unproved placeholder that
introduces \texttt{sorryAx}.

Absence of such placeholders is necessary, but is not a complete review of a
formal result. One must still read the theorem's type. A theorem of the form
``if a uniform gap exists, then a uniform gap exists'' can compile without
any placeholder and yet supply no new mathematics. An empty index type can
make a supposedly universal statement vacuous. An unconnected construction
witness can prove an existential object while failing to connect it to the
intended lattice approximants. These are specification issues, not failures
of the kernel. The external Yang--Mills review discussed in the main text
illustrates why keeping the construction witness in the acceptance target
matters.

The present file makes each assumption explicit with
\texttt{set\_option autoImplicit false}. Its declarations are short enough
to inspect directly. This improves auditability, but it does not erase the
need to compare their types with the intended mathematical prose.

\section{Rebuilding and extending the companion}

The tracked \texttt{lean-toolchain} and \texttt{lakefile.toml} specify the
compiler and exact mathlib revision. A conventional build obtains that
toolchain, resolves the pinned package, obtains the corresponding compiled
cache if desired, and runs \texttt{lake build}. The README also records the
actual environment used for this verification. Downloaded toolchains and
mathlib caches are dependencies, not vendored proof evidence.

Any extension should add a theorem-to-text mapping and preserve the axiom
audit. For example, formalizing the nodal periodic resolution would require
a quotient algebra, a basis equivalence and module maps, followed by kernel
and image computations. Formalizing the torus-link topology would require
continuous maps and an inverse homeomorphism. Formalizing the Fourier
sampling theorem would require finite-dimensional complex Hilbert spaces,
a normalized state, a unitary operator and a probability interpretation.
Each is a meaningful project. Replacing those requirements by a single
axiom bearing the desired theorem's name would not satisfy them.

\section{Seven additional quantum and binary-chain declarations}

The additional file \path{QuantumInterfaces.lean} was compiled against the
same pinned toolchain. Its seven declarations have printed theorem types
and foundational dependency reports in \texttt{quantum-\allowbreak verification.log};
the independent \texttt{quantum-\allowbreak verification.json} records its hash and
inventory. The original 25-declaration file and its verification manifest
remain unchanged. These two files supply the first 32 selected declarations;
the second-pass module below adds sixteen.

Three theorems verify $X^2=Z^2=I$ and $XZ=-ZX$ for the actual integer
matrices $X=\left(\begin{smallmatrix}0&1\\1&0\end{smallmatrix}\right)$ and
$Z=\left(\begin{smallmatrix}1&0\\0&-1\end{smallmatrix}\right)$. Further
declarations define the sign character $\chi$ of $\mathbb F_2$ and prove
\[
 \chi(a+b)=\chi(a)\chi(b).
\]
For a binary vector space $C$, a linear functional $\ell:C\to\mathbb F_2$,
and an amplitude function $\psi:C\to R$ with $R$ any commutative ring, the
implemented operators are
\[
 (X_a\psi)(v)=\psi(v+a),\qquad
 (Z_\ell\psi)(v)=\chi(\ell(v))\psi(v).
\]
The theorem \path{shift_phase_commutation} derives
$X_aZ_\ell=\chi(\ell(a))Z_\ell X_a$ as equality of functions. The general
coefficient ring includes complex amplitudes, but the formal declaration
does not assert an additional analytical Hilbert-space theorem.

For actual maps $d_2:C_2\to C_1$ and $d_1:C_1\to C_0$ over $\mathbb F_2$
with $d_1d_2=0$, \path{boundary_phase_shift_commute} sets $a=d_2f$ and
uses the functional $\ell d_1$. Their pairing vanishes by the given
composition equation, so the corresponding shift and phase operators
commute. The result thus verifies a concrete algebraic interface used in
CSS constructions; it does not assume the desired commutation as its own
premise. The seventh theorem derives
\[
 E-(V-1)-(F-1)=2g
\]
from the separately supplied integer equation $V-E+F=2-2g$. Surface
classification, the ranks of cellular boundary maps, and error-correction
performance are outside this arithmetic substitution. No complete quantum
circuit or QNN training procedure is claimed formalized.

\section{Sixteen additional diagram and transport declarations}

The second-pass file \path{DiagramGrammar.lean} was compiled with exit
status zero under the same pinned Lean 4.24.0 and mathlib environment.
Its actual output is \path{diagram-verification.log}; its declaration and
axiom inventory is \path{diagram-verification.json}. The source hash is
\begin{center}\scriptsize
\texttt{871d34de90d99a263094f7abc4aadb51f5012d613c671b04fff43f5a0c35c1c6}.
\end{center}
All 48 declarations across the three modules were also freshly rechecked
during the second pass, with source hashes preserved in the audit ledger.

The typed syntax has objects, identity maps, typed generators, and sequential
composition. A context has one typed hole. The declarations
\path{context_sound}, \path{rewriteStep_sound}, and
\path{derivation_sound} prove that semantically equal rules remain equal in
such contexts and through finite derivations. Structural identity and
associativity rules have their own soundness proof. An explicit family of
linear-equivalence generators adds inverse cancellation, and
\path{equiv_derivation_sound} verifies finite derivations using those actual
rules. These statements verify the sequential component of
Theorem~\ref{thm:typed-rewrite-soundness}; they do not encode the full graphical
tensor, dagger, trace, measurement, or channel syntax.

Separate finite-matrix declarations verify tensor interchange, trace
factorization, rectangular trace cyclicity, and the dimension value of a
closed identity loop. They are algebraic component theorems, not an implicit
coherence theorem for arbitrary drawings. The complex-transport declaration
uses actual linear equivalences to preserve a zero composite. The
exactness-transport declaration preserves image--kernel equality, and
\path{homotopy_implies_exact} derives exactness from an explicit contracting
homotopy identity. These hypotheses are meaningful data, not names for the
desired conclusion.

Finally \path{matrix_hodge_conjugacy} proves a finite complex-matrix identity
with the necessary endpoint conditions $U_0^\dagger U_0=I$ and
$U_2^\dagger U_2=I$. Its algebraic identity allows an arbitrary middle
matrix $U_1$; unitarity of $U_1$ is the further condition used for the
spectral interpretation in Proposition~\ref{thm:unitary-chain-transport}.
The formal statement therefore supports the displayed finite conjugacy
component. It does not assert an infinite-dimensional gap or a continuum
Yang--Mills result.

\section{Complete verified companion sources}

The following listing reproduces the verified file so that the mathematical
types, proof scripts and audit commands can be read together. The separate
source file is the authoritative build input.

\lstinputlisting{../lean/BridgeTheorems.lean}

The additional quantum and binary-chain file follows. Its source and audit
are independent of the preceding file.

\lstinputlisting{../lean/QuantumInterfaces.lean}

The second-pass typed diagram, matrix contraction, and complex-transport
module follows. Its exact statement types and foundational dependency
commands are part of the authoritative source.

\lstinputlisting{../lean/DiagramGrammar.lean}
